\section{Post-Study Existing-Record Analyses}
\label{app:derived}
The major outcomes were known before these calculations; the earlier descriptors below were added in v0.8.6. Before recomputation, a written plan fixes the complete retained inputs, formulas, ordering and treatment of missing records. These are secondary descriptive or sensitivity analyses; no original endpoint, primary CI or bootstrap stream is replaced. Anonymous \texttt{derived/inputs.json}, \texttt{derived/results.json} and \texttt{analysis.py} retain the earlier analyses; \texttt{verify\_v090.py} independently recomputes the new extents and budgets from the same retained records.


\subsection{Complete-trajectory extent and reconvergence}
These descriptions use all original arrays after their aggregate outcomes are known. For B570, each of 24 prompts has three stable repeats with 128-ID outputs. Prediction equals live in every run. On the four reference-divergent prompts, every position from the first difference through cap differs from reference:
\begin{table}[t]
\caption{B570 non-reference forecast extent on four distinct trajectories. Indices are zero-based; counts include the first different ID. Repeats confirm stability, not additional prompt units.}
\label{tab:b570extent}
\centering
\begin{tabular}{@{}lrrr@{}}
\toprule
Task & First index & Extent through cap & Later IDs \\
\midrule
/134 & 88 & 40 & 39 \\
/42 & 73 & 55 & 54 \\
/145 & 103 & 25 & 24 \\
/138 & 93 & 35 & 34 \\
\midrule
Total & --- & \DiagBNonref & \DiagBAfter \\
\bottomrule
\end{tabular}
\end{table}
Token positions are correlated within these capped trajectories. This extent does not establish an uncapped endpoint or a numerical cause.

For RTX V4 A, /25's UB2 and live arrays each have 75 IDs; they differ only at index 4 and then reconverge. On /15, reference and UB2 are identical 28-ID arrays; both disagree with live on 10 positions of its 25-ID extent and emit three additional IDs. On /18, reference has 72 IDs versus 69 for UB2/live, with 40 differences on the common extent and three additional IDs. H2/R\_ONLY matches live in all cases. The main table shows these extents next to the original counts. A hindsight choice between the two already measured complete target-only arrays covers 18/19 prompts but still fails on /15. Choosing with knowledge of the live answer is not a deployable predictor, and the comparison does not exhaust target-only shape/history policies.

\subsection{B570 raw logits, margins and continuation}
Each reserved run retains three raw target-logit rows for a common proposal-conditioned prefix, whether or not all proposals are reached by actual acceptance. Top-one/top-two are ranked over the full vocabulary before sampler transformations; margin is their score difference. A fixed proposal's raw logit is also retained (absent for the bonus row). These logits are neither normalized probabilities nor sampler-adjusted margins. Numerical fields change in all 72 runs; timestamps/metadata are excluded. The 216 rows and 195 reachable paired decisions are correlated, as are three repeats per prompt.
\begin{table}[t]
\caption{B570 post-study raw-logit descriptors over all 216 paired rows. $\Delta$ = UB128 minus UB1; margins use retained top-two scores/IDs. Units are raw logits.}
\label{tab:scorederived}
\centering
\begin{tabular}{lrrr}
\toprule
Statistic & Min & Median & Max \\
\midrule
$\Delta$ top-one & -0.4275 & -0.0443 & 2.3503 \\
$|\Delta|$ top-one & 0.0019 & 0.0695 & 2.3503 \\
UB1 margin & 0.2291 & 17.5021 & 33.6660 \\
UB128 margin & 0.2003 & 17.5128 & 33.5351 \\
\bottomrule
\end{tabular}
\end{table}
Top-two ID changes occur in three rows; raw top-one IDs never change. Absolute top-one delta quartiles are 0.0388 and 0.1327, with p95 0.3437. Positive minimum raw margins coexist with unchanged immediate decisions; these summaries do not prove a whole-generation null effect or locate a numerical operation.

The new cohort is HumanEval/1, /8, /13, /21, /29, /36, /41, /51, /71, /74, /75, /79, /83, /84, /85, /92, /94, /110, /124, /131, /152, /153, /155, /163. Every primary and alternate endpoint reaches cap 128, with no earlier EOG. Source-before, restored and source-after serialized KV content hashes/lengths match at all 72 forks; consumed lengths and positions match. Recursive sampler cloning is source-audited, with matching recorded chain/seed/history. Six same-partition UB1-to-UB1 R0 restores have equal scores and full continuations; the complete R0 set has 24 existing executions including original/native equivalence. Sampler opaque-state bytes are not separately serialized. These controls support the sampled contrast, not identity of unavailable external binaries.

\begin{table}[t]
\caption{All six B570 continuation-sensitive prompts. Zero-based generated positions, distance = first difference minus fork index. Every row is stable across all three repetitions, with 128/128-ID capped endpoints. Length-only divergence would occur at the shorter array's length; none occurs here.}
\label{tab:forkpositions}
\centering
\begin{tabular}{lrrr}
\toprule
HumanEval & Fork & First difference & Distance \\
\midrule
/29 & 65 & 114 & 49 \\
/75 & 65 & 80 & 15 \\
/84 & 65 & 115 & 50 \\
/110 & 65 & 84 & 19 \\
/131 & 64 & 74 & 10 \\
/163 & 66 & 97 & 31 \\
\bottomrule
\end{tabular}
\end{table}
All sampled immediate decisions agree. Later UB128 target blocks persist and proposals evolve with the alternate accepted history. These secondary differences concern another cohort; they cannot explain old H1's four divergences or a single-block restoration.

\subsection{M4 opportunity between two fixed modes}
From all 540 raw held-out observations, the three-condition medians exactly reproduce the original 60 prompt rows. Target-only is faster than Always on 2/60 prompts, both Spec-Bench; there are no equal Target/Always medians. Always/oracle GM is 1.0006919671$\times$, where oracle is the per-prompt lower Target or Always median. Its implied GM latency reduction from Always is $1-1/1.0006919671=0.0006914886$, about \DiagMOpportunity\%. This compares only the two measured fixed modes and is not an arbitrary switching-policy bound.
\begin{table}[t]
\caption{The complete M4 target-faster subset, from original medians. Ratios are relative latencies ($\times$); lower is better. All six Guard runs fall back at round 8 but each Guard median remains above target-only. No transition component is inferred.}
\label{tab:m4opportunity}
\centering
\begin{tabular}{lrrr}
\toprule
Task & Target (s) & Always/Target & Guard/Target \\
\midrule
/559 & 11.9085 & 1.0270 & 1.0215 \\
/95 & 10.9203 & 1.0150 & 1.0029 \\
\bottomrule
\end{tabular}
\end{table}
Guard/Always paired prompt ratios have min/median/max 0.9706/1.0022/1.3297 and quartiles 0.9985/1.0260. Family GMs are HumanEval 1.0099 and Spec-Bench 1.0446; these 30/30 families are descriptive, not stable small-stratum estimates.

A separate sensitivity stream uses Python Random seed 2026100101, 10,000 paired draws, sampling each family's 30 prompts with replacement in HumanEval then Spec-Bench order. Each draw computes GM directly; linear .025/.975 percentiles give Guard/Always [1.0135, 1.0423] and Always/oracle [1.0000, 1.0018]. Guard remains slower; original pooled primary and standalone secondary intervals stay unchanged. No CI is formed by dividing interval endpoints.

\subsection{358H conditional zero-draft accounting}
For each path, sum matched baseline time $B$, complete target work $T$, draft $D_t$, residual $X_t$, committed IDs $N$ and rounds $L$ over all 90 processed observations. $T_b=B/N$, $C=N/L$, and $R_{\mathrm{target}}=T/(LT_b)$. At the same progress, target work and baseline, with nonnegative draft/residual,
\[
S=\frac{C}{D_{\mathrm{real}}+R_{\mathrm{target}}+X}
\le\frac{C}{R_{\mathrm{target}}}=\frac{B}{T}.
\]
Both paths retain 11,520 committed IDs; rounds are 4,677 CPU and 4,686 NPU. Full-precision processed times yield bounds CPU 0.8274529834 and NPU 0.8351858524$\times$. Even removing both draft and residual costs cannot reach parity under these fixed conditions. This is arithmetic total-time accounting, not a bound on prompt-GM speedup, a measured optimization, separable replay saving, persistence gain or a bound after changing target execution/proposals.


\subsection{358H original-prompt-estimator ceiling and budget}
For prompt $i$ and repetition $j$, let $B_{ij}$ be matched baseline time, $T_{ij}$ the complete target region, and $D_{ij}+X_{ij}\ge0$ the draft/residual regions. With the measured baseline, target work and output progress fixed, any counterfactual retaining nonnegative other cost has $S'_{ij}\le B_{ij}/T_{ij}$. Median and GM are monotone on positive values, so the original estimator satisfies
\[
G'=\operatorname{GM}_i\operatorname{median}_j S'_{ij}
\le U=\operatorname{GM}_i\operatorname{median}_j(B_{ij}/T_{ij}).
\]
All 180 unique processed rows (30 prompts $\times$ three per path) are retained. Positive baseline/target time, nonnegative draft/residual, and component sums within 0.01ms are checked. The canonical processed CSV and the earlier anonymous export agree field for field. Full-precision $U$ is 0.8372192092432174 CPU and 0.8448401086171792 NPU. The distinct pooled-total bounds above remain unchanged.

Neither path has a prompt-median ceiling at or above one (0/30 each). CPU's maximum is 0.9736196764; NPU's is 0.9999089575 on /64. One individual NPU repetition has $B/T=1.000163695$, so the result does not assert all 90 NPU run-level ceilings below one or a robust margin for every prompt. These are conditional arithmetic on one session, not confidence intervals or a population forecast.

If target work uniformly falls to $(1-\delta)T_{ij}$ with other cost zero, the prompt-GM ceiling becomes $U/(1-\delta)$. Reaching one therefore needs $\delta\ge1-U$: \DiagCPUReduction\% CPU and \DiagNPUReduction\% NPU. Actual remaining overhead can require more. Under the different pooled-total estimator the thresholds are 17.2547\% and 16.4814\%; the units must not be mixed. Changing target execution, proposal quality or committed progress invalidates the fixed-work premise. No persistence saving, target-work reduction or new qualified optimization is measured.

\subsection{Earlier RTX two-path absolute latency (30 prompts)}
All 180 original condition times are milliseconds; medians retain three repetitions per prompt and the original 30-prompt pairing. The following distributions summarize 30 condition medians, not 180 independent prompts.
\begin{table}[t]
\caption{Earlier two-path RTX cost study: post-study absolute prompt-median times (ms), 30 prompts. The original target/serialized GM and CI remain unchanged; the ratio of these distribution medians is not that GM.}
\label{tab:rtxabsolutetime}
\centering
\begin{tabular}{lrrrr}
\toprule
Condition & Min & Q25 & Median & Q75 \\
\midrule
Target & 295.2 & 1186.9 & 1907.9 & 2730.3 \\
Serialized & 560.3 & 1821.1 & 2923.3 & 4053.6 \\
\bottomrule
\end{tabular}
\end{table}
Maxima are 2939.7 and 4490.6 ms. The JSON output also reports all 90 individual times per condition separately. Independent primary replay is 0.6458157742698285; the original report's one-ulp 0.6458157742698286 and original CI/seed remain recorded. The post-prefill clock boundary and measured hook probe in Section N apply; 0.6 $\mu$s p99 is about 0.0012\% of the shortest 49.9572-ms R0 interval. That is a clock-call probe, not all logging overhead. The complete native-log files have median line counts 3516.5 target and 4353.5 serialized, and median byte lengths 266908 and 339335. Committed-ID event counts match for all 90 pairs. File-wide counts include initialization outside the timer; they do not measure charged logging cost or imply cancellation. No production-speed estimate is made.
